\SourcePage{74}{77}
\section{Fundamental properties of the Schrödinger equation}

The conditions to be satisfied by solutions of the Schrödinger equation are
quite general in nature. The wave function, to begin with, must be
single-valued and continuous throughout%
\SourcePage{75}{78}
\relax\ all of space. The continuity requirement remains in force even when the
field $U(x,y,z)$ itself has surfaces of discontinuity. On any such surface,
the wave function together with its derivatives must still be continuous. The latter are
not continuous, however, if the potential energy $U$ becomes infinite beyond
some surface. A particle cannot enter at all a region of space where
$U=\infty$, that is, $\psi=0$ must hold everywhere in that region. Continuity
of $\psi$ requires it to vanish on the boundary of this region; in this case,
however, the derivatives of $\psi$ generally undergo a jump.

If the field $U(x,y,z)$ nowhere becomes infinite, the wave function must also
be finite throughout space. The same condition must be observed when $U$
becomes infinite at some point, provided it does not do so too rapidly, namely
as $1/r^s$ with $s<2$ (see also \secref{35}).

Let $U_{\min}$ be the minimum value of the function $U(x,y,z)$. Since the
particle's Hamiltonian is the sum of two terms, the kinetic- and
potential-energy operators $\hat T$ and $U$, the mean energy in an arbitrary
state is the sum $\bar E=\bar T+\bar U$. But all eigenvalues of the operator
$\hat T$ (which coincides with the Hamiltonian of a free particle) are
positive, and therefore the mean value $\bar T>0$ as well. Taking account also
of the evident inequality $\bar U>U_{\min}$, we find that
$\bar E>U_{\min}$. Since this inequality holds for every state, it plainly
holds for all energy eigenvalues as well:
\begin{equation}
  E_n>U_{\min}.
  \tag{18.1}
\end{equation}

Suppose a particle moves through a force field that goes to zero at large
distances; following standard convention, we fix $U(x,y,z)$ so that it too
tends to zero at infinity. One readily verifies that the
eigenvalue spectrum for negative energies is then discrete: every state with
$E<0$ in a field that vanishes at infinity is a bound state. Indeed, in stationary states of the continuous spectrum,
which correspond to infinite motion, the particle is at infinity (see
\secref{10}). At sufficiently large distances the presence of the field may be
neglected and the particle's motion may be regarded as free; in free motion,
however, the energy can only be positive.

Conversely, the positive eigenvalues form a continuous spectrum and
correspond to infinite motion; for $E>0$ the Schrödinger equation generally
does not have%
\SourcePage{76}{79}
\relax\ (in the field under consideration) solutions for which the integral
$\int|\psi|^2\,dV$ converges\footnote{From a purely mathematical standpoint,
however, it must be qualified that, for certain forms of the function
$U(x,y,z)$ (having no physical significance), a discrete set of values may
drop out of the continuous spectrum.}.

We draw attention to the fact that, in quantum mechanics, a particle in finite
motion may also be present in regions of space where $E<U$. Although the
probability $|\psi|^2$ of finding the particle tends rapidly to zero deeper
inside such a region, it nevertheless differs from zero at every finite
distance. This represents a fundamental departure from classical
mechanics, where a particle is altogether unable to penetrate a region with $U>E$. In
classical mechanics, penetration into this region is impossible because, for
$E<U$, the kinetic energy would become negative, meaning that the velocity
would be imaginary. In quantum mechanics, too, the kinetic-energy eigenvalues are
positive; nevertheless, no contradiction arises here, because if the
measurement process localizes the particle at some definite point of space,
that same process disturbs the particle's state in such a way that it ceases
altogether to have any definite kinetic energy.

If $U(x,y,z)>0$ throughout space (with $U\to0$ at infinity), inequality (18.1)
gives $E_n>0$. On the other hand, since the spectrum must be continuous for
$E>0$, it follows that in the case under consideration the discrete spectrum
is absent altogether, that is, only infinite motion of the particle is
possible.

Suppose that at some point, chosen as the coordinate origin, $U$ tends to
$-\infty$ according to the law
\begin{equation}
  U\approx-\alpha/r^s,
  \qquad \alpha>0.
  \tag{18.2}
\end{equation}
Consider a wave function that is finite in a certain small region of radius
$r_0$ around the origin and zero outside it. The uncertainty in the particle's
coordinate values in such a wave packet is of order $r_0$; the uncertainty in
the momentum value is therefore $\sim\hbar/r_0$. The mean kinetic energy in
this state is of order $\hbar^2/mr_0^2$, and the mean potential energy is of
order $-\alpha/r_0^s$. Suppose first that $s>2$. Then, for sufficiently small
$r_0$, the sum $\hbar^2/mr_0^2-\alpha/r_0^s$ takes negative values of
arbitrarily large absolute magnitude. But if the mean energy can take such
values, this means in any event that there exist negative energy eigenvalues%
\SourcePage{77}{80}
\relax\ of arbitrarily large absolute magnitude. Energy levels with large
$|E|$ correspond to motion of the particle
in a very small region of space around the origin. The ``normal'' state will
correspond to the particle being located at the origin itself, that is, the
particle will ``fall'' into the point $r=0$.

If $s<2$, on the other hand, the energy cannot take negative values of
arbitrarily large absolute magnitude. The diagonal spectrum starts from a certain
finite negative value, and the particle does not fall to the center.
We note that in classical mechanics a particle can in principle fall to the
center in any attractive field (that is, for any positive $s$). The case
$s=2$ will be considered separately in \secref{35}.

Next, let us investigate the character of the energy spectrum as a function of
the behaviour of the field at large distances. Suppose that as
$r\to\infty$ the potential energy, while negative, tends to zero according to
the power law (18.2) (in this formula $r$ is now large). Consider a wave packet
``filling'' a spherical shell of large radius $r_0$ and thickness
$\Delta r\ll r_0$. Then again the order of magnitude of the kinetic energy is
$\hbar^2/m{\left(\Delta r\right)}^2$, and that of the potential energy is
$-\alpha/r_0^s$. Let $r_0$ increase while $\Delta r$ is increased at the same
time (so that $\Delta r$ grows in proportion to $r_0$). If $s<2$, then for
sufficiently large $r_0$ the sum
$\hbar^2/m{\left(\Delta r\right)}^2-\alpha/r_0^s$ becomes negative. It follows
that stationary states with negative energy exist in which the particle may
be found with appreciable probability at large distances from the origin. But
this means that negative energy levels of arbitrarily small absolute magnitude
exist (it must be remembered that wave functions decay rapidly in the region
of space where $U>E$). Thus, in the case under consideration, the discrete
spectrum contains an infinite set of levels which crowd together towards the
level $E=0$.

If, however, the field falls off at infinity as $-1/r^s$ with $s>2$, there
are no negative levels of arbitrarily small absolute magnitude. The discrete
spectrum ends at a level with nonzero absolute magnitude, so that the total
number of levels is finite.

The Schrödinger equation for the stationary-state wave functions $\psi$, together
with the conditions placed on its solutions, is real. Its solutions can
therefore always be chosen real\footnote{This is not true for systems in a
magnetic field.}. The eigenfunctions of nonde%
\SourcePageJoin{78}{81}
generate energy values automatically prove to be real up to an immaterial
phase factor. Indeed, $\psi^*$ satisfies the same equation as $\psi$ and is
therefore also an eigenfunction for the same value of the energy; hence, if
that value is nondegenerate, $\psi$ and $\psi^*$ must be essentially
identical, that is, they may differ at most by a constant factor whose modulus
equals unity. Wave functions that correspond to one and the same degenerate level need
not be real, but a set of real functions can always be obtained by a suitable
choice of their linear combinations.

The full (time-dependent) wave functions $\Psi$, however, are determined by an
equation whose coefficients contain $i$. This equation nevertheless retains
its form if $t$ is replaced by $-t$ and complex conjugation is performed at
the same time\footnote{It is assumed that the potential energy $U$ has no
explicit dependence on time: the system is either closed or is in a constant
(nonmagnetic) field.}. The functions $\Psi$ can therefore always be chosen so
that $\Psi$ and $\Psi^*$ differ only in the sign of time.

As is known, the equations of classical mechanics are unchanged under
\emph{time reversal}, i.e.\ upon reversing the sign of time. In quantum
mechanics, as we see, symmetry with respect to the two directions of time
manifests itself in the invariance of the wave equation under reversal of
the sign of $t$ accompanied by the replacement of $\Psi$ with $\Psi^*$. It must be
remembered, however, that here this symmetry applies only to the equations and
not to the concept of measurement itself, which has a fundamental role in
quantum mechanics (as discussed in detail in \secref{7}).
